\documentclass[11pt]{article}
\usepackage[margin=1in]{geometry}
\usepackage{amsmath,amssymb,amsthm}
\usepackage{graphicx}
\usepackage{booktabs}
\usepackage{xcolor}
\usepackage{listings}
\usepackage[numbers]{natbib}
\usepackage{hyperref}

\newtheorem{theorem}{Theorem}
\newtheorem{lemma}[theorem]{Lemma}
\theoremstyle{definition}
\newtheorem{definition}[theorem]{Definition}
\theoremstyle{remark}
\newtheorem*{remark}{Remark}

\emergencystretch=2em

\lstdefinelanguage{lean}{
  morekeywords={field_simp,ring,div_nonneg,def,theorem,rcases,haveI,subst,nlinarith,linarith,have,refine,intro,induction,with,simp,only,rw,unfold,exact,example,namespace,end,open,fun,Prop,by,noncomputable,section,import},
  sensitive=true,
  morecomment=[l]{--},
  morecomment=[s]{/-}{-/},
}
\lstset{
  language=lean,
  basicstyle=\ttfamily\footnotesize,
  keywordstyle=\color{blue!60!black}\bfseries,
  commentstyle=\color{gray},
  columns=fullflexible,
  keepspaces=true,
  breaklines=true,
  frame=single,
  extendedchars=true,
  literate=
    {ℝ}{{$\mathbb{R}$}}1
    {ℕ}{{$\mathbb{N}$}}1
    {→}{{$\to$}}1
    {↔}{{$\leftrightarrow$}}1
    {∀}{{$\forall$}}1
    {∃}{{$\exists$}}1
    {∈}{{$\in$}}1
    {∑}{{$\sum$}}1
    {∏}{{$\prod$}}1
    {∧}{{$\wedge$}}1
    {≤}{{$\le$}}1
    {≥}{{$\ge$}}1
    {≠}{{$\ne$}}1
    {⨆}{{$\bigsqcup$}}1
    {∉}{{$\notin$}}1
    {←}{{$\leftarrow$}}1
    {ρ}{{$\rho$}}1
    {⟨}{{$\langle$}}1
    {⟩}{{$\rangle$}}1
    {·}{{$\cdot$}}1
    {𝓝}{{$\mathcal{N}$}}1
    {₀}{{$_0$}}1
    {⊢}{{$\vdash$}}1
}

\title{A Recovery-Adjusted Work Floor:\\ A Lean Proof that Recovered Supply Lowers, and Exactly Attains, the Fresh-Supply Requirement}
\author{Justin Hart, Viridis LLC}
\date{17 November 2026}

\begin{document}
\maketitle

\begin{abstract}
A process must deliver a useful-work target $T$ from two sources: a \emph{fresh} supply $w_f$
and a \emph{recovered} supply $w_r$, converted at modeled efficiencies $e_f$ and $e_r$ that share
a common cap $e_{\max}>0$. We give a Lean~4 proof, against Mathlib, of the sealed statement
\texttt{P3.RecoveryWorkFloor}. Part~1 (floor): if the useful work $e_fw_f+e_rw_r$ reaches $T$,
then $w_f\ge (T-e_{\max}w_r)/e_{\max}$. Part~2 (attainment): whenever $e_{\max}w_r\le T$, the
fresh supply $(T-e_{\max}w_r)/e_{\max}$ is non-negative and, run at the cap, delivers exactly $T$,
so the floor cannot be raised. The proof derives the floor from a restatement of
\texttt{target\_requires\_supply} from the Useful and Dissipated Work Ledger (Run-163): total supply
must be at least $T/e_{\max}$, and the recovered part moves across the inequality. The floor is a
bound on \emph{modeled} efficiency only. It is not a Landauer bound, it does not use the
Intelligence Bound conjecture, and it makes no claim about any real efficiency or recovered supply.
\end{abstract}

\section{Introduction}
Recovering material or energy from residues reduces what must be drawn from fresh sources. That
is the operational core of the circular economy \citep{geissdoerfer2017} and of industrial
symbiosis, where one firm's residue becomes another's input \citep{chertow2000}. A decision
kernel that credits recovery needs a precise answer to a simple question: \emph{if a target amount
of useful work must be delivered and some of the supply is recovered, how much fresh supply is
still required at minimum?}

Under a linear conversion model with a common efficiency cap the answer is a one-line formula,
\[
w_f\ \ge\ \frac{T-e_{\max}w_r}{e_{\max}}=\frac{T}{e_{\max}}-w_r ,
\]
and the bound is tight. Each unit of recovered supply lowers the fresh-supply floor by exactly one
unit, and no more. The formula is elementary. What a kernel needs is the exact list of hypotheses
under which it may be applied automatically, and the assurance that the floor is the best possible
one. Project P3 of Viridis Research proposed \texttt{P3.RecoveryWorkFloor} in session P3-T001
(5~November 2026), and the program owner sealed it under decision D-013. This note gives a Lean
proof that the sealed statement holds exactly as written.

\paragraph{Contributions.}
\begin{itemize}
\item A Lean~4 proof of \texttt{ViridisStatements.P3.RecoveryWorkFloor.Statement}
(Section~\ref{sec:proof}), structured as an extension of the Run-163 lemma
\texttt{target\_requires\_supply} \citep{hart2026udwl}. It depends only on the standard axioms
\texttt{propext}, \texttt{Classical.choice} and \texttt{Quot.sound}.
\item A short human-readable proof (Section~\ref{sec:math}) that records which hypotheses each part
actually uses.
\item A plain statement of what kind of bound this is and is not (Section~\ref{sec:limits}).
\end{itemize}

\section{Background and related work}
\paragraph{Circular economy and symbiosis.}
\citet{geissdoerfer2017} review the circular economy as a sustainability paradigm, centred on
substituting recovered resources for virgin ones. \citet{chertow2000} defined industrial symbiosis
as the exchange of materials, energy, water and by-products among firms. Earlier P3 sessions gave
Lean proofs of two material-accounting statements. \texttt{P3.LoopMassBalance} (P3-T002) says that
with recovery fraction $r$ the virgin input is $D(1-r)$. \texttt{P3.CascadedRecovery} (P3-T003)
says that the virgin share of $n$ stages in series is $\prod_i(1-r_i)$, which is no larger than
any single stage's share. This note turns from material to work.

\paragraph{The Useful and Dissipated Work Ledger.}
Run-163 of the Viridis canon \citep{hart2026udwl} models two supplies $w_1,w_2\ge0$ converted at
efficiencies $e_1,e_2$, with useful work $e_1w_1+e_2w_2$ and modeled dissipation
$(1-e_1)w_1+(1-e_2)w_2$. Its lemma \texttt{target\_requires\_supply} says that if both efficiencies
are at most $e_{\max}>0$ and useful work reaches $T$, then the total supply is at least
$T/e_{\max}$. \texttt{P3.RecoveryWorkFloor} reads $w_1$ as fresh and $w_2$ as recovered supply,
and asks for the bound on the fresh part alone, together with its tightness.

\paragraph{What this is not.}
\citet{landauer1961} showed that logically irreversible operations entail a minimum heat
generation of order $kT$ per operation. That is a physical lower bound that holds for every
device. The floor proved here is a different kind of statement. Its only physics is the modeling
assumption $e\le e_{\max}$, where $e_{\max}$ is a number chosen by the modeler. It is not a
Landauer bound. It also does not rely on the program's Intelligence Bound, which remains a
conjecture.

\paragraph{Formal tools.}
The proof is written in Lean~4 \citep{moura2021lean4} (toolchain v4.28.0) against Mathlib
\citep{mathlib2020} (commit \texttt{8f9d9cff}).

\section{The model and the statement}
\label{sec:model}
\begin{definition}[Recovery work model]
For fresh supply $w_f$, recovered supply $w_r$ and efficiencies $e_f,e_r$, the \emph{useful work}
is $U(w_f,w_r,e_f,e_r)=e_fw_f+e_rw_r$. For a target $T$, recovered supply $w_r$ and cap
$e_{\max}$, the \emph{fresh floor} is $F(T,w_r,e_{\max})=(T-e_{\max}w_r)/e_{\max}$.
\end{definition}

\begin{theorem}[\texttt{P3.RecoveryWorkFloor}]\label{thm:main}
\begin{enumerate}
\item (Floor.) Let $w_f,w_r\ge0$, $0\le e_f\le e_{\max}$, $0\le e_r\le e_{\max}$ and $e_{\max}>0$.
If $T\le U(w_f,w_r,e_f,e_r)$, then $F(T,w_r,e_{\max})\le w_f$.
\item (Attainment.) Let $e_{\max}>0$, $w_r\ge0$ and $e_{\max}w_r\le T$. Then
$F(T,w_r,e_{\max})\ge0$ and $U(F(T,w_r,e_{\max}),w_r,e_{\max},e_{\max})=T$.
\end{enumerate}
\end{theorem}

\begin{lstlisting}[caption={Definitions and \texttt{Statement} from the sealed module \texttt{ViridisStatements.P3.RecoveryWorkFloor} (docstrings omitted).}]
def useful (wf wr ef er : ℝ) : ℝ := ef * wf + er * wr

def freshFloor (T wr emax : ℝ) : ℝ := (T - emax * wr) / emax

def Statement : Prop :=
  (∀ wf wr ef er emax T : ℝ, 0 ≤ wf → 0 ≤ wr →
    0 ≤ ef → ef ≤ emax → 0 ≤ er → er ≤ emax → 0 < emax →
    T ≤ useful wf wr ef er → freshFloor T wr emax ≤ wf) ∧
  (∀ wr emax T : ℝ, 0 < emax → 0 ≤ wr → emax * wr ≤ T →
    0 ≤ freshFloor T wr emax ∧ useful (freshFloor T wr emax) wr emax emax = T)
\end{lstlisting}

\section{Mathematical proof}
\label{sec:math}
\begin{lemma}[Run-163, \texttt{uniform\_efficiency\_bound}]\label{lem:ueb}
If $w_f,w_r\ge0$ and $e_f,e_r\le e_{\max}$, then $U(w_f,w_r,e_f,e_r)\le e_{\max}(w_f+w_r)$.
\end{lemma}
\begin{proof}
$e_{\max}(w_f+w_r)-U=(e_{\max}-e_f)w_f+(e_{\max}-e_r)w_r$, a sum of products of non-negative
numbers.
\end{proof}

\begin{lemma}[Run-163, \texttt{target\_requires\_supply}]\label{lem:trs}
Under the hypotheses of Lemma~\ref{lem:ueb}, if also $e_{\max}>0$ and $T\le U$, then
$T/e_{\max}\le w_f+w_r$.
\end{lemma}
\begin{proof}
$T\le U\le e_{\max}(w_f+w_r)$; divide by $e_{\max}>0$.
\end{proof}

\begin{proof}[Proof of Theorem~\ref{thm:main}]
(1) For $e_{\max}\ne0$, $F(T,w_r,e_{\max})=T/e_{\max}-w_r$. By Lemma~\ref{lem:trs},
$T/e_{\max}-w_r\le w_f$.

(2) $T-e_{\max}w_r\ge0$ and $e_{\max}>0$, so $F\ge0$. At the cap,
$U(F,w_r,e_{\max},e_{\max})=e_{\max}\cdot\frac{T-e_{\max}w_r}{e_{\max}}+e_{\max}w_r=T$.
\end{proof}

\begin{remark}[Which hypotheses are used]
Part~(1) does not use $e_f\ge0$ or $e_r\ge0$; the sealed statement includes them because they are
part of the physical model. It does use $w_r\ge0$. If $w_r<0$ and $e_r<e_{\max}$, the term
$(e_{\max}-e_r)w_r$ is negative and the bound can fail. Part~(2) does not use $w_r\ge0$; it uses
$e_{\max}w_r\le T$, and the sealed module's negative control shows that this condition is needed:
with $e_{\max}=1/2$, $w_r=4$ and $T=1$ the floor is $-2$. The cap is also essential. The sealed
module shows that with efficiencies $1$ above a cap of $1/2$, the target $8$ is reached with fresh
supply $4$, below the floor $12$.
\end{remark}

\paragraph{Interpretation.} With cap $e_{\max}=1/2$ and target $T=8$, no recovery needs at least
$16$ units of fresh supply. Four units of recovered supply lower the floor to $12$, and twelve
units at the cap reach $8$ exactly, while eleven reach only $15/2$. Every unit of recovered supply
reduces the fresh-supply floor by exactly one unit. That is so because both sources share one cap.
Recovery is credited one for one, no more, and only up to the point $w_r=T/e_{\max}$ where recovery
alone covers the target.

\section{The Lean development}
\label{sec:proof}
The candidate file \texttt{proofs/P3.RecoveryWorkFloor.lean} follows Section~\ref{sec:math}. The
proof may import only Mathlib and the sealed module, so the Run-163 file cannot be imported. Its
two lemmas \texttt{uniform\_efficiency\_bound} and \texttt{target\_requires\_supply} are restated
over the sealed \texttt{useful} and proved the same way: \texttt{nlinarith} with the two gap terms,
then \texttt{div\_le\_iff\textsubscript{0}}. \texttt{freshFloor\_eq} rewrites the floor as
$T/e_{\max}-w_r$ by \texttt{field\_simp}. Part~1 is then \texttt{linarith}. Part~2 is
\texttt{div\_nonneg} followed by \texttt{field\_simp} and \texttt{ring}.

\begin{lstlisting}[caption={The complete candidate proof \texttt{proofs/P3.RecoveryWorkFloor.lean}.}]
import Mathlib
import ViridisStatements.P3.RecoveryWorkFloor

namespace Candidate

open ViridisStatements.P3.RecoveryWorkFloor

/-- Uniform efficiency bound (Run-163 `uniform_efficiency_bound`, restated for the sealed
`useful`): useful work is at most `emax · (wf + wr)`. -/
theorem uniform_efficiency_bound (wf wr ef er emax : ℝ) (hwf : 0 ≤ wf) (hwr : 0 ≤ wr)
    (hef : ef ≤ emax) (her : er ≤ emax) :
    useful wf wr ef er ≤ emax * (wf + wr) := by
  have h1 : 0 ≤ (emax - ef) * wf := mul_nonneg (sub_nonneg.mpr hef) hwf
  have h2 : 0 ≤ (emax - er) * wr := mul_nonneg (sub_nonneg.mpr her) hwr
  unfold useful
  nlinarith

/-- Run-163 `target_requires_supply`, restated: reaching `T` needs total supply `≥ T / emax`. -/
theorem target_requires_supply (wf wr ef er emax T : ℝ) (hwf : 0 ≤ wf) (hwr : 0 ≤ wr)
    (hef : ef ≤ emax) (her : er ≤ emax) (he : 0 < emax) (hT : T ≤ useful wf wr ef er) :
    T / emax ≤ wf + wr := by
  have hu := uniform_efficiency_bound wf wr ef er emax hwf hwr hef her
  rw [div_le_iff₀ he]
  linarith [mul_comm emax (wf + wr)]

/-- The floor is the total-supply bound with the recovered part moved across:
`(T - emax·wr)/emax = T/emax - wr`. -/
theorem freshFloor_eq (T wr emax : ℝ) (he : 0 < emax) :
    freshFloor T wr emax = T / emax - wr := by
  unfold freshFloor
  field_simp

/-- Part 1 (floor). -/
theorem floor (wf wr ef er emax T : ℝ) (hwf : 0 ≤ wf) (hwr : 0 ≤ wr)
    (hef : ef ≤ emax) (her : er ≤ emax) (he : 0 < emax)
    (hT : T ≤ useful wf wr ef er) : freshFloor T wr emax ≤ wf := by
  rw [freshFloor_eq T wr emax he]
  linarith [target_requires_supply wf wr ef er emax T hwf hwr hef her he hT]

/-- Part 2 (attainment): the floor is nonnegative and, at `ef = er = emax`, delivers `T`. -/
theorem attain (wr emax T : ℝ) (he : 0 < emax) (hT : emax * wr ≤ T) :
    0 ≤ freshFloor T wr emax ∧ useful (freshFloor T wr emax) wr emax emax = T := by
  unfold freshFloor useful
  refine ⟨div_nonneg (by linarith) he.le, ?_⟩
  field_simp
  ring

theorem main : ViridisStatements.P3.RecoveryWorkFloor.Statement :=
  ⟨fun wf wr ef er emax T hwf hwr _ hef _ her he hT =>
      floor wf wr ef er emax T hwf hwr hef her he hT,
   fun wr emax T he _ hT => attain wr emax T he hT⟩

end Candidate
\end{lstlisting}

\section{Results}
\label{sec:results}
We ran the program's checker on the candidate file. It runs the same stages as the verifier:
denylist, compile, statement match, axioms and kernel replay. The output was
\texttt{PASSED (preview)}, and the reported axioms were \texttt{propext},
\texttt{Classical.choice} and \texttt{Quot.sound}. A first, direct version of the proof (without
the Run-163 lemmas) also passed on its first attempt. We then restructured it so that the floor is
visibly a corollary of \texttt{target\_requires\_supply}, and it passed again. No automated prover
was used. The certification page that precedes this paper gives the verifier's own judgement.

\section{What this settles and what it does not}
\label{sec:limits}
\paragraph{Settled.} Assuming the verifier accepts the file, the floor and its attainment hold
exactly as sealed, for all real inputs that satisfy the hypotheses.

\paragraph{Not settled.}
\begin{itemize}
\item \textbf{A modeled-efficiency bound.} $e_{\max}$ is an input chosen by the modeler. The theorem
says what follows if every conversion efficiency is at most $e_{\max}$. It does not say what
$e_{\max}$ is for any real process. It is not a Landauer bound or any other physical limit
\citep{landauer1961}. It does not follow from, and gives no evidence for, the Intelligence Bound
conjecture.
\item \textbf{No empirical content.} Nothing here estimates any real efficiency, recovered supply
or target. Those stay assumed.
\item \textbf{One common cap.} Recovered work is often of lower quality than fresh work, so a
per-source cap $e_r\le e_{r,\max}<e_{f,\max}$ would be more realistic. Then the floor becomes
$(T-e_{r,\max}w_r)/e_{f,\max}$. That variant is not sealed and is not proved here.
\item \textbf{Recovery is free in the model.} The model does not charge any work for recovering
$w_r$. Coupling $w_r$ to the loop and cascade models (for example, $w_r$ proportional to
recovered mass $rD$) and charging recovery work is left for later work.
\end{itemize}

\section{Reproducibility}
\sloppy
The proof file \texttt{proofs/P3.RecoveryWorkFloor.lean} imports only \texttt{Mathlib} and the
sealed module \texttt{ViridisStatements.P3.RecoveryWorkFloor}. It compiles under toolchain
\texttt{leanprover/lean4:v4.28.0} and the Mathlib commit below.
\begin{lstlisting}[language={},frame=none]
sealed module sha256:
  5f5a53efa57459cfd1bd276db7c9ceb96886c8d3f54a66b87b907311f7a7252c
Mathlib commit:
  8f9d9cff6bd728b17a24e163c9402775d9e6a365
viridis-lean-check proofs/P3.RecoveryWorkFloor.lean P3.RecoveryWorkFloor
\end{lstlisting}

\section{Changelog}
\begin{itemize}
\item 2026-11-05 (P3-T001): first version. Proposed P3.LoopMassBalance, P3.CascadedRecovery,
P3.RecoveryWorkFloor and P3.CyclicSymbiosis as statement modules, with controls, probe results
and rationale.
\item 2026-11-15 (P3-T002): P3.LoopMassBalance sealed (D-011). Narrowed the paper to a note on
this statement with a full Lean proof and a human-readable proof.
\item 2026-11-16 (P3-T003): P3.CascadedRecovery sealed (D-012). Note on the cascade statement with
a full Lean proof.
\item 2026-11-17 (P3-T004): P3.RecoveryWorkFloor sealed (D-013). This version is a note on the
recovery-adjusted work floor with a full Lean proof built on Run-163's
\texttt{target\_requires\_supply}. Added references to Run-163 and to Landauer (1961). Earlier
results are summarised in the background section; their proofs stay in the earlier run packages.
\end{itemize}

\bibliographystyle{plainnat}
\bibliography{refs}
\end{document}
